\subsubsection{Overzicht en classificatie van detectiekenmerken}
\label{subsubsec:Overzicht en classificatie van detectiekenmerken}
Automatische phishingdetectie in e-mails is de voorbije jaren sterk geëvolueerd. Vroege detectiesystemen steunden op handmatig gedefinieerde kenmerken die als invoer dienden voor machinelearningmodellen. Zo introduceerden \textcite{Fette2007} PILFER, een classificatiemodel dat op basis van tien handmatig gekozen kenmerken bepaalt  of een e-mail phishing is. \textcite{Bergholz2010} bouwden hierop voort en combineerden  inhoudelijke en structurele kenmerken, zoals statistische modellen voor de onderwerpen van een e-mail en de sequentiële analyse van tekst en externe links. Later onderzoek maakt  daarnaast gebruik van machine- en deep-learningmodellen \autocite{Salloum2022}.